\clearpage
\fancyhead[L]{\small SONNENPARALLAXE.}
\SourceLeaf{528}
\begin{center}\small Göttingische gelehrte Anzeigen. Stück 17. Seite 161..166. 1810 Januar 29.\\[3pt]\rule{.72\linewidth}{.3pt}\end{center}

Hr \textsc{Lars Regner}, Professor der Astronomie zu Upsala, Mitglied der dortigen königl. Societät der Wissenschaften, hat unter der Ueberschrift: \emph{Supplementum ad historiam de Parallaxeos Solaris inventione, illustriss.\ scientiarum Societatis regiae Gottingensis dijudicationi subjectum}, einen Aufsatz eingesandt, wovon Folgendes im Wesentlichen der Inhalt ist. Hr. \textsc{Regner} erzählt zuerst, er sei schon im Jahre 1806 von der Unrichtigkeit der von den Astronomen angenommenen Sonnenparallaxe überzeugt gewesen, und habe im Anfange des folgenden Jahres eine kleine Schrift darüber an die Petersburger Academie der Wissenschaften eingesandt. Hr. \textsc{Fuss}, als Secretär dieser Academie, habe geantwortet, dass die Abhandlung Hrn. \textsc{Schubert} zur Beurtheilung übergeben und von diesem als irrig und falsch befunden sei; er selbst zweifle auch gar nicht, dass die blosse Andeutung der Quelle der Fehler hinlänglich sein werde, Hrn. \textsc{Regner} davon zu überführen; überdies erbiete er sich zur Mittheilung des ganzen Berichts selbst. Diesen wartete aber Hr. \textsc{Regner} nicht ab, sondern liess, in der Ueberzeugung, dass der angedeutete Tadel ungerecht sei, seine Schrift mit einigen Erläuterungen in Upsala drucken, und sandte sie dann von neuem nach Petersburg. Hierauf schickte ihm Hr. \textsc{Fuss} ohne Weiteres den Bericht des Hrn. \textsc{Schubert} selbst.

Hr. \textsc{Regner} lässt nun einen Auszug aus seiner (uns sonst nicht bekannt gewordenen) Druckschrift folgen. Das Wesentliche davon ist die Bemerkung, dass, wenn man die Parallaxe der Sonne zwischen $8''5$ und $8''9$ voraussetzt, die Attraction der Sonne auf den Mond beträchtlich grösser ist, als die Attraction der Erde auf denselben. Aus dieser (allerdings richtigen) Bemerkung glaubt Hr. \textsc{Regner} schliessen zu müssen, dass jene Parallaxe unmöglich richtig sein könne, weil der Mond von der Sonne stärker, als von der Erde, angezogen (nach Hrn. \textsc{Regner}'s Vorstellung), diese verlassen, und seine eigene Bahn um die Sonne beschreiben müsste. Ja, wenn man auch nur annähme, dass der Mond von der Sonne nur eben so stark angezogen würde, wie von der Erde, so würden im Neumonde beide Anziehungen einander vollkommen aufheben, der Mond nach der Tangente fortgehen, und dadurch sofort in eine überwiegende Attractionssphäre der Sonne gerathen. Diese letztere Behauptung will Hr. \textsc{Regner} durch eine Figur deutlich machen, wo er die Erde, den Mond und die Sonne in einer geraden Linie durch $T$, $b$, $S$ bezeichnet, und dann zeigt, dass die Punkte der in $b$ errichteten senkrechten geraden Linie sich von $T$ schneller entfernten, als von $S$. (Wie konnte Hr. Prof. \textsc{Regner} vergessen, dass die Erde, während der Mond in der Tangente fortginge, ja nicht in $T$ bleibt, sondern vermöge ihrer Tangential-Geschwindigkeit von selbst nachrückt, also keinesweges der Abstand des Mondes von ihr nothwendig abzunehmen braucht? Eben so liegt in einer \emph{grössern} Anziehung nach der Sonne zu nichts Widersprechendes. Es folgt \emph{blos}, dass die Mondbahn im absoluten Raume zur Zeit des Neumondes gegen die Sonne zu concav werden muss; und dass sie dies wirklich ist, davon kann man sich leicht überzeugen, wenn man sie auch nur mit Vernachlässigung der Ungleichheiten der Bewegung der Erde und des Mondes als eine Epicycloide betrachtet.)

Hr. \textsc{Regner} fügt jetzt einen Auszug aus Hrn. \textsc{Schubert}'s Bericht bei. Hr. \textsc{Regner} müsse, urtheilt \textsc{Schubert}, die Abhandlung entweder in grosser Eile oder aus Scherz geschrieben haben; denn wäre es ernstlich damit gemeint, so würde folgen, dass der Verfasser weder die Methode, wodurch die Sonnenparallaxe bestimmt worden ist, noch die \emph{ersten Grundsätze} des Problems der drei Körper kenne. Die erstere Behauptung rechtfertigt \textsc{Schubert} damit, dass man in den Beobachtungen des Venusdurchganges einen Fehler von 27 Minuten annehmen muss, wenn die Sonnenparallaxe 21 Secunden wäre, welche Hr. \SourceLeaf{529}\textsc{Regner} seiner vorhin erwähnten Vorstellung nach, für die kleinste zulässige hält. Zur Bekräftigung der zweiten Behauptung zeigt \textsc{Schubert}, dass Hr. \textsc{Regner} die absolute Bewegung des Mondes mit der relativen um die Erde verwechselt habe, indem es, wenn man letztere betrachtet, nicht auf die Attraction des Mondes von der Sonne, sondern lediglich darauf ankömmt, wie viel diese anders ist, als die Attraction der Erde von der Sonne: dieser Unterschied aber gegen die Attraction des Mondes von der Erde nur gering sei. (Eine weitere Auseinandersetzung dieser \textsc{Schubert}'schen Entwickelung werden Leser, die in den Anfangsgründen der physischen Astronomie nicht ganz fremd sind, uns gern erlassen.)

Hierauf lässt nun endlich Hr. \textsc{Regner} folgen: Animadversiones in relationem D\textsuperscript{nı} \textsc{Schubert} ad Scient. Imperial. Petropolitanam de Scripto parallaxin Solis tractante. Wir übergehen die langen Klagen über Ungerechtigkeit, wodurch er sich von Hrn. \textsc{Schubert} gekränkt und beschimpft glaubt, um nur noch kürzlich dasjenige anzuzeigen, womit Hr. \textsc{Regner} die Beschuldigungen von sich abzuwälzen sucht. Er bemüht sich zuvörderst, die Genauigkeit der Beobachtungen bei den Durchgängen der Venus 1761 und 1769, welche \textsc{Lalande} auf ein paar Secunden in Zeit, oder auf ein paar Zehntheile einer Secunde in Bogen, gesetzt hatte, verdächtig zu machen; eine solche Genauigkeit sei mit den damals angewandten Fernröhren schlechterdings unerreichbar gewesen, und \emph{deswegen} müsse man eine Anhäufung von Beobachtungsfehlern bis auf 27 Zeitminuten ganz wohl für möglich halten, und die damals gemachten Beobachtungen zur Bestimmung der Sonnenparallaxe für ganz untauglich erklären! (Wir können uns unmöglich entschliessen, über eine solche Art, sich zu rechtfertigen, Etwas hinzu zu setzen.)

Hr. \textsc{Regner} glaubt nun, dass ihm nichts weiter übrig sei, als noch Hrn. \textsc{Schubert}'s zweite Beschuldigung zu entkräften. Er fängt mit einigen Gemeinplätzen über die Nachtheile der blinden Unterwerfung unter Auctorität und hergebrachte Irrthümer in den Wissenschaften an, zu denen, seiner Meinung nach, eben hier der Ort sei. Nemlich blos auf \textsc{Newton}'s Wort hin glaube man seitdem, dass die Wirkung der Attraction der Sonne auf die Erde und auf den Mond in gleichen Abständen gleich gross sei; allein dieser Satz sei ganz falsch, und Hr. \textsc{Regner} setzt an dessen Stelle folgenden: Wenn zwei Körper von ungleichen Massen durch gleiche Kräfte sollicitirt werden, so bekommen sie ungleiche und den Massen verkehrt proportionale Geschwindigkeiten. (Warum muss man dem Hrn. Professor \textsc{Regner} hier in Erinnerung bringen, dass dieser Satz nur dann wahr sein kann, wenn von gleicher \emph{bewegender} Kraft die Rede ist, wie man sie in der Mechanik von beschleunigender Kraft unterscheidet; und dass die Sonne auf die Weltkörper nur in so fern bewegende Kraft ausübt, als sie auf deren einzelne Theile wirkt, also desto mehr bewegende Kraft auf das Ganze, je mehr Theile angezogen werden?) Hr. \textsc{Regner} läugnet, um diesen Satz zu vertheidigen, geradezu die Beweiskraft der Versuche, die mit fallenden Körpern unter der Glocke der Luftpumpe gemacht worden, weil theils die Fallzeiten zu kurz seien, theils doch nie die Luft \emph{ganz} herausgeschafft werden könne. Dass unter der Glocke die Feder eben so geschwind fällt, als der Ducaten, kömmt, nach Hrn. \textsc{Regner}'s Meinung, blos daher, weil die noch übrige Luft der Feder mehr Widerstand entgegen setze, als dem Goldstück, und im völlig leeren Raume würde die Feder so viel schneller fallen, als sie weniger Masse habe. Eben so glaubt er, das Phänomen, dass zwei Kugeln von gleichem Volumen und ungleichen Massen, an gleich langen Fäden aufgehängt, gleich schnelle Pendelschwingungen machen (was übrigens, in so fern der Widerstand der Luft noch merklich ist, gar nicht wahr ist), beweise gerade seinen Satz unwidersprechlich, weil die kleinere Masse von einer so viel grössern Kraft sollicitirt sein müsse, um denselben Widerstand eben so leicht zu überwinden und eben so schnell zu fallen. Zuletzt findet Hr. \textsc{Regner} noch einen Beweis seines Satzes in dem Gleichgewichte zweier Körper von ungleichen Massen an ungleichen und diesen umgekehrt proportionalen Armen des Hebels, weil die Virtualgeschwindigkeiten hier im umgekehrten Verhältnisse der Massen stehen. Indem nun Hr. \textsc{Regner} durch dies alles es als ausgemacht ansieht, dass \SourceLeaf{530}die Wirkung der Attraction der Sonne auf den Mond und auf die Erde bei jenem um so viel grösser ist, als bei dieser, als die Masse des erstern kleiner ist, als die der andern, findet er es leicht, die Unzulässigkeit von Hrn. \textsc{Schubert}s Schlüssen zu zeigen.

Weder der Raum, noch die Bestimmung dieser Blätter erlauben uns, Etwas weiter als diese Darlegung zu geben, und unsern Lesern in den Folgerungen, die sie leicht selbst daraus ziehen werden, vorzugreifen.
